\cotefile{157-2}
\scanpage{323}
\mspage{62}
% archive sheet number: 322
% ms. 1–2
de type $\Delta_+$, ou aussi au \deleted{\unsure{n'importe quel}}\inserted{chemin dégénéré} de n'importe quel type $\tau$, disons
% ms. 3
\diagram{\begin{tikzcd}[ampersand replacement=\&,column sep=small]
x \& x\arrow[l,"\mathrm{id}"']\arrow[r,"\mathrm{id}"] \& x \arrow[r,dash,"\mathrm{id}"] \& x \& x\arrow[l,"\mathrm{id}"'] \& x\arrow[l,"\mathrm{id}"']{,}
\end{tikzcd}}
% ms. 4
un tel chemin n'étant autre que $p_\tau^*(x)$, où
% ms. 5
\[p_\tau:\tau\longrightarrow 1\]
% ms. 6–7
est le morphisme de \unsure{dégénérescence} (l'unique morphisme dans $\underline{\unsure{\mathrm{Ty}|\mathrm{Tz}}}$ vers $1$, qui est objet final),
% ms. 8–9
d'où
\diagram{\begin{tikzcd}[ampersand replacement=\&,row sep=0pt]
p_\tau^*:\underline{\mathrm{Ch}}_1(\deleted{\unsure{x}}x,x)=\{x\}\arrow[r] \& \underline{\mathrm{Ch}}_\tau(x,x)\\
\& \cap\\
\& \underline{\mathrm{Ch}}(x,x).
\end{tikzcd}}

% ms. 10–12
La catégorie \inserted{$\Pi(X)$} ainsi obtenue doit s'appeler le \uline{groupoïde fondamental} de $X$. En effet, \deleted{\unsure{a)}}

% ms. 12–15
a) toutes les flèches de $\Pi(X)$ sont inversibles, i.e. $\Pi(X)$ est bien un groupoïde, et

% ms. 16
\unsure{b)} le foncteur canonique
% ms. 17
\[X\longrightarrow\Pi(X)\]
% ms. 18–22
qui est l'identité sur les objets, et qui est défini sur les flèches par $\unsure{u\mapsto\mathrm{cl}\,u|u\mapsto\mathrm{cl}(u)}$ comme ci-dessus, est \uline{universel} \inserted{(dans $\mathrm{Cat}$)} pour les foncteurs \inserted{(i.e. les $\mathrm{Cat}$-morphismes)} de $X$ vers des \uline{groupoïdes}.
% ms. 23–26
(Plus généralement, il est universel pour les foncteurs $X\xrightarrow{F}Y$ tels que $F$ transforme \unsure{tte} flèche de $X$ en flèche inversible.)
